\chapter{Step 06: The Serial Line}
\label{ch:step06}

% Listings with line numbers, for listings whose lines the text refers to.
\lstdefinestyle{numbered}{numbers=left,numberstyle=\tiny\color{codecomment},numbersep=6pt}

So far the board has talked to us through LEDs and a display, and we have
talked to it through buttons and switches. To load programs, we need a real
connection to the computer. Step 06 uses the serial line that the USB chip
provides (Section~\ref{sec:board}). The design of \cc{06_io_mirror} receives
every character we type in a terminal, shows its code on the display, and
sends it back. So what we see in the terminal is what the FPGA sent. On the
way we meet a lot of new SystemVerilog: modules with \cc{parameter}s, state
machines with \cc{typedef enum}, a \cc{function}, and a trap that the
synthesis tool does not let us fall into.

\section{Sending Bits One After the Other}
\label{sec:step06-uart}

The serial line has two wires, one for each direction. In
\cc{icebreaker.pcf} they are called \cc{RX} (pin 6, from the computer to the
FPGA) and \cc{TX} (pin 9, from the FPGA to the computer), named from the
FPGA's point of view. The circuit that sends and receives on such a line is
called a \term{UART}, a \emph{universal asynchronous receiver/transmitter}.

\emph{Asynchronous} means: there is no clock wire. Sender and receiver only
agree on the \term{baud rate}, the number of bits per second. Ours is
9600\,baud, so one bit lasts $1/9600\,\mathrm{s} \approx 104.2\,\mu\mathrm{s}$,
which is $12\,000\,000 / 9600 = 1250$ cycles of our 12\,MHz clock. Each side
measures this time with its own clock.

One byte is sent as a \term{frame} (Figure~\ref{fig:step06-frame}):
\begin{itemize}
\item When nothing is sent, the line is \emph{idle} and stays at 1.
\item A \term{start bit}, a 0, announces a frame. Its falling edge is the
  moment from which the receiver measures time.
\item Then come the 8 data bits, the \emph{least} significant bit first.
\item A \term{stop bit}, a 1, ends the frame. Our transmitter sends two of
  them; a receiver needs only one, the second simply looks like idle line.
\end{itemize}
A frame with one stop bit takes 10 bit times, ours 11. The receiver samples
each bit in its middle, where it is furthest from both of its edges. Since the
two clocks are never exactly equal, the sample points drift; but if the baud
rates of both sides differ by a few percent at most, the last sample of a
frame, about $9.5$ bit times after the falling edge, is still well inside its
bit. The next frame starts with a new falling edge, so the error does not
accumulate from frame to frame.

\begin{figure}[ht]
\centering
\begin{tikzpicture}[x=9.5mm,y=5mm]
  % bits of 'A' = 0x41, LSB first
  \def\lo{0} \def\hi{1}
  \draw[black!15] (0,0) -- (13,0);
  \draw[signal] (0,1) -- (1,1) -- (1,0) -- (2,0) -- (2,1) -- (3,1) -- (3,0)
     -- (8,0) -- (8,1) -- (9,1) -- (9,0) -- (10,0) -- (10,1) -- (13,1);
  \foreach \x/\lab in {1/start,2/$d_0$,3/$d_1$,4/$d_2$,5/$d_3$,6/$d_4$,7/$d_5$,8/$d_6$,9/$d_7$,10/stop,11/stop} {
    \draw[black!25] (\x,-0.4) -- (\x,1.6);
    \node[font=\scriptsize] at (\x+0.5,1.9) {\lab};
  }
  \draw[black!25] (12,-0.4) -- (12,1.6);
  \foreach \x/\v in {1/0,2/1,3/0,4/0,5/0,6/0,7/0,8/1,9/0,10/1,11/1} {
    \node[font=\scriptsize\ttfamily] at (\x+0.5,-0.8) {\v};
  }
  \foreach \x in {1,...,10} { \fill[green!45!black] (\x+0.5,0.5) circle (1.2pt); }
  \node[font=\scriptsize,anchor=east] at (0,0.5) {\cc{TX}};
  \node[note] at (0.5,1.9) {idle};
  \node[note] at (12.5,1.9) {idle};
  \draw[decorate,decoration={brace,mirror,amplitude=3pt}] (1,-1.4) -- (2,-1.4)
     node[midway,below=2pt,note] {1250 cycles};
\end{tikzpicture}
\caption{The frame for the character \cc{A}, code $\mathtt{41}_{16} =
\mathtt{01000001}_2$, as \cc{uart_tx} sends it: start bit, the data bits
$d_0$ to $d_7$ starting with the least significant, two stop bits. The dots
mark where a receiver samples.}
\label{fig:step06-frame}
\end{figure}

\section{The Top Module}
\label{sec:step06-top}

The top module \cc{test.sv} connects a receiver \cc{uart_rx}, a transmitter
\cc{uart_tx} and \cc{dev_hex}, and controls them with a small state machine.
Its first part declares the signals between them:

\svpart{style=numbered,firstline=1,firstnumber=1,lastline=33}{06_io_mirror/test.sv}

The comments \cc{/* verilator lint_off UNUSEDSIGNAL */} and \cc{lint_on} tell
Verilator not to warn that \cc{rx_idle} and \cc{rx_eop} are never read. The
receiver provides them, and we do not need them.

\textbf{\cc{typedef enum} names the states.} Lines 26 to 31 define a new type
\cc{state_t} whose values have names: \cc{IO_READY_TO_RECEIVE},
\cc{IO_BUFFERING_DATA}, and so on. The names stand for the numbers 0, 1, 2,
3, in this order. Line 33 declares two signals of this type, \cc{state} with
the power-up value \cc{IO_READY_TO_RECEIVE} and \cc{next_state}. As with
\cc{logic}, the type says nothing about the hardware; that \cc{state} becomes a
register follows from the \cc{always_ff} block that assigns it. What the
\cc{enum} gives us is readability: we write the name of a state instead of a
number, and a simulator shows the names in the waveform.

How many bits does \cc{state} have? Without further information, an
\cc{enum} is based on the type \cc{int}, which has 32 bits. sv2v writes
\cc{reg [31:0] state}. yosys finds that the upper 30 bits are always 0 and
removes their flip-flops, but we will see in
Section~\ref{sec:step06-latch} that the 32 bits still matter.

\subsection*{A state machine}

A \term{state machine} is a register, the \emph{state}, together with
combinational logic that computes the next state from the current state and
the inputs. We have seen the smallest example in Chapter~\ref{ch:step02}.
Here the description is split into three blocks:

\svpart{style=numbered,firstline=35,firstnumber=35,lastline=71}{06_io_mirror/test.sv}

The \cc{always_comb} block says which state follows which: a combinational
circuit from \cc{state} to \cc{next_state}, here simply the successor in a
cycle of four. The \cc{always_ff} block is the state register. In some states
it takes over \cc{next_state} only under a condition, and if the condition is
false it keeps its value. As a circuit, this is a register with enable, where
the enable is 1 in \cc{IO_READY_TO_RECEIVE} if \cc{rx_ready} is 1, in
\cc{IO_BUFFERING_DATA} if \cc{tx_busy} is 0, and in the other two states
always. Figure~\ref{fig:step06-fsm} shows the result as a state diagram.

\begin{figure}[ht]
\centering
\begin{tikzpicture}[node distance=12mm and 44mm,
    st/.style={box,fill=blue!7,minimum width=38mm,minimum height=8mm,font=\footnotesize}]
  \node[st] (r) {\cc{IO_READY_TO_RECEIVE}};
  \node[st,right=of r] (b) {\cc{IO_BUFFERING_DATA}};
  \node[st,below=of b] (i) {\cc{IO_INIT_TRANSMISSION}};
  \node[st,below=of r] (s) {\cc{IO_STARTED_TRANSMISSION}};
  \draw[arrow] (r) -- node[above,note]{\cc{rx_ready}} node[below,note]{\cc{data_buf <= rx_data_out}} (b);
  \draw[arrow] (b) -- node[right,note]{\cc{!tx_busy}} (i);
  \draw[arrow] (i) -- node[above,note]{always} node[below,note]{\cc{tx_start <= 1}} (s);
  \draw[arrow] (s) -- node[left,note]{always}  node[right,note]{\cc{tx_start <= 0}} (r);
\end{tikzpicture}
\caption{The state machine of \cc{06_io_mirror/test.sv}. An arrow is taken at a
rising edge of \cc{CLK} if its condition holds; otherwise the state stays. The
assignments below the arrows are made at the same edge.}
\label{fig:step06-fsm}
\end{figure}

The third block describes what happens in each state:

\svpart{style=numbered,firstline=73,firstnumber=73,lastline=94}{06_io_mirror/test.sv}

It describes three registers with enable: \cc{data_buf} loads the received
byte when \cc{rx_ready} is 1 in the first state; \cc{tx_data_in} and
\cc{tx_start} change in the last two states. A state that is not mentioned,
or the empty statement \cc{;}, means: the register keeps its value.

Together: the receiver signals a byte with \cc{rx_ready}, which is 1 for
exactly one clock cycle. At that edge, the byte goes into \cc{data_buf}, and
the machine moves to \cc{IO_BUFFERING_DATA}. There it waits until the
transmitter is free. In \cc{IO_INIT_TRANSMISSION} it hands the byte to the
transmitter and sets \cc{tx_start}, which is 1 during the one cycle in
\cc{IO_STARTED_TRANSMISSION} and is cleared at the edge back to the first
state. So the transmitter sees a start pulse of one cycle. The display shows
\cc{data_buf}, the last byte received.

\begin{exercise}
\label{ex:step06-next}
Change one line of \cc{test.sv} so that the board answers each character with
the next one in the code table: \cc{a} with \cc{b}, \cc{A} with \cc{B}.
\end{exercise}

\section{Parameters}
\label{sec:step06-parameter}

The last part of \cc{test.sv} creates the receiver and the transmitter:

\svpart{style=numbered,firstline=96,firstnumber=96,lastline=120}{06_io_mirror/test.sv}

Between the module name and the instance name stands a list
\texttt{\#(\,\dots)}: values for the \term{parameters} of the module. In
\cc{uart_rx.sv} they are declared as
\begin{lstlisting}
parameter clk_freq = 12000000;
parameter baud = 115200;
parameter oversampling = 8;
\end{lstlisting}
A \cc{parameter} is a constant of the module, like a \cc{localparam}
(Chapter~\ref{ch:step03}), with one difference: each instance may replace its
value. The numbers after \cc{=} are only defaults. Here \cc{.baud(BAUD)} sets
\cc{baud} to 9600 for both instances, \cc{oversampling} keeps its default.
Parameters can also be given by position, in the order of their declaration;
\cc{uart_rx} does that when it creates its tick generator:
\begin{lstlisting}
baud_tick_gen #(clk_freq, baud, oversampling) tickgen(.clk(clk), .enable(1'b1), .tick(os_tick));
\end{lstlisting}

A parameter is not a wire. It is a number that is fixed when the design is
synthesised, and it changes the circuit itself: other values give other
widths and other constants, so each set of parameter values describes a
different circuit. yosys builds a separate module for each set, with a name
derived from the values. The per-module statistics of \cc{make report}
(Appendix~\ref{app:tools}) list them:

\outputfile{step06_paramod.txt}

There are two versions of \cc{baud_tick_gen}, one with oversampling 8 for the
receiver and one with oversampling 1 for the transmitter.

\textbf{A note on style.} \cc{uart_rx.sv}, \cc{uart_tx.sv} and \cc{baud_tick_gen.sv} come from the
iCEBreaker examples of Piotr Esden-Tempski and are written in the older
Verilog style. \cc{reg} and \cc{wire} are the predecessors of \cc{logic}, and
\cc{always @(posedge clk)} with \cc{<=} is what we write as \cc{always_ff}, as
in \cc{dev_hex} (Chapter~\ref{ch:step03}). Ports without a type, such as
\cc{input clk}, are single wires. It describes the same kind of hardware; only
the tools no longer check that an \cc{always} block really describes
flip-flops.

\section{Baud Ticks from \texorpdfstring{12\,MHz}{12 MHz}}
\label{sec:step06-baud}

Both the transmitter and the receiver need a signal that is 1 for one clock
cycle, at the baud rate (transmitter) or at eight times the baud rate
(receiver). For 9600\,baud a counter to 1250 would do, but for eight times
9600 we would need $12\,000\,000 / 76\,800 = 156.25$ cycles, and for 115200\,baud
$104.17$. \cc{baud_tick_gen} solves this for all rates at once:

\svpart{style=numbered,firstline=23,firstnumber=23,lastline=48}{06_io_mirror/baud_tick_gen.sv}

\textbf{The circuit} is in lines 39, 45 and 46: a register \cc{acc} of $w + 1$
bits, where $w$ is \cc{acc_width}, and an adder. In each cycle, the low $w$
bits of \cc{acc} plus the constant \cc{inc} are stored back; the top bit
\cc{acc[w]} receives the carry-out of this $w$-bit addition, and it is the
output \cc{tick}. Mathematically, with $a_n$ the value of the low $w$ bits
after $n$ cycles,
\[
  a_{n+1} = (a_n + \mathit{inc}) \bmod 2^w, \qquad
  \mathit{tick}_{n+1} = 1 \iff a_n + \mathit{inc} \ge 2^w .
\]
So $a_n = n \cdot \mathit{inc} \bmod 2^w$ (starting from 0), and the number of
ticks in the first $n$ cycles is $\lfloor n \cdot \mathit{inc} / 2^w \rfloor$.
At a clock of $f$\,Hz this gives on average
\[
  f \cdot \frac{\mathit{inc}}{2^w}
\]
ticks per second. Since $\mathit{inc} < 2^w$, at most one tick comes per
cycle, and the gap between two ticks is either $\lfloor 2^w /
\mathit{inc}\rfloor$ or $\lceil 2^w / \mathit{inc}\rceil$ cycles. Such a
circuit is called a \term{phase accumulator}, or a fractional divider: it
divides the clock by the fraction $2^w / \mathit{inc}$.

\textbf{The constants.} To get $o \cdot b$ ticks per second for baud rate $b$
and oversampling $o$, we need $\mathit{inc} \approx o \cdot b \cdot 2^w / f$,
rounded. The larger $w$, the smaller the rounding error. Line 37 sets $w$ to 8
more than the number of binary digits of $f / b$. Line 42 computes the rounded
value of $o \cdot b \cdot 2^w / f$ in integer arithmetic: adding half the
divisor, \cc{clk_freq >> (shiftlimiter + 1)}, before the division turns
truncation into rounding, and \cc{shiftlimiter} shifts numerator and
denominator so that the intermediate values stay below $2^{31}$. If we
uncomment the \cc{$display} in line 43, the simulation prints the values for
our two instances:

\outputfile{step06_baud_inc.txt}

For the transmitter, $w = 19$ and $\mathit{inc} = 419$, so
\[
  12\,000\,000 \cdot \frac{419}{2^{19}} \approx 9590.1\ \text{ticks per second},
\]
about 0.1\,\% slower than 9600; a bit lasts 1251 or 1252 cycles instead of 1250.
For the receiver, $\mathit{inc} = 3355$ gives $76\,789.9$ instead of $76\,800$,
an error of $0.013\,\%$. Both are far inside the tolerance of a UART.

\textbf{A function computed by the tools.} \cc{log2_} in lines 31 to 34 is a
\cc{function}: it takes an integer and returns one. Here it is only used in
\cc{localparam}s, so it is evaluated by the tools, during synthesis, and it
describes \emph{no} circuit at all. Its \cc{while} loop is a loop in yosys,
not in hardware. (A function applied to signals would describe combinational
logic; we will not need that.) Despite its name, \cc{log2_(v)} is not
$\log_2 v$: it counts the binary digits of $v$, that is, $\lfloor \log_2 v
\rfloor + 1$ for $v \ge 1$. For example, $f/b = 1250$ has 11 binary digits,
so $w = 19$.

\textbf{The enable.} If \cc{enable} is 0, line 45 sets \cc{acc} to \cc{inc}.
The transmitter connects \cc{enable} to \cc{tx_busy}, so its tick generator
starts afresh with every frame, and the first tick comes one bit time after
the transmission starts. The receiver's tick generator always runs.

\begin{exercise}
\label{ex:step06-baud}
With the default parameters (12\,MHz, 115200\,baud, no oversampling), compute
\cc{acc_width}, \cc{shiftlimiter} and \cc{inc}, and the actual tick rate. How
large is the relative error?
\end{exercise}

\section{The Transmitter}
\label{sec:step06-tx}

\svpart{style=numbered,firstline=27,firstnumber=27,lastline=49}{06_io_mirror/uart_tx.sv}

The states of the transmitter are not an \cc{enum} but \cc{localparam}s with
codes chosen by hand. The choice is clever: the eight data states have bit 3
set and the number of the data bit in bits 2 to 0, the states in which the
line is 1 (\cc{IDLE}, \cc{BIT_STOP1}, \cc{BIT_STOP2}) are the codes below 4,
and \cc{BIT_START} is 4. This makes the output logic small:

\svfile[style=numbered,firstline=51,firstnumber=51,lastline=82]{06_io_mirror/uart_tx.sv}

\cc{tx_shift} is a \term{shift register}: when the transmitter is idle and
\cc{tx_start} is 1, it loads the byte; in each data state it shifts right by
one at every \cc{bit_tick}. So \cc{tx_shift[0]} is always the bit to be sent,
first $d_0$, then $d_1$, and so on: least significant bit first. The state
moves on at every \cc{bit_tick}. In line 82, \cc{tx_state < 4} is 1 in idle
and in the stop bits, \cc{tx_state[3] & tx_shift[0]} is the data bit in the
data states, and in \cc{BIT_START} both are 0. The codes 1, 5, 6 and 7 are
not used, and the \cc{default} branch leads from them back to \cc{IDLE}.

\section{The Receiver}
\label{sec:step06-rx}

Receiving is harder: the receiver has to find the start of a frame and the
middle of each bit by itself. It does that by looking at the line eight
times per bit, at each \cc{os_tick} of its tick generator. This is called
\term{oversampling}.

\svpart{style=numbered,firstline=55,firstnumber=55,lastline=80}{06_io_mirror/uart_rx.sv}

\textbf{Synchronising.} \cc{RX} changes whenever the computer sends, without
any relation to our clock. A flip-flop whose input changes just at the clock
edge may need unusually long to settle on 0 or 1. The two flip-flops
\cc{rx_sync} in series give it a whole sampling period to settle before the
value is used. This is the standard way to bring an outside signal into a
clocked design.

\textbf{Filtering.} \cc{filter_cnt} is a 2-bit counter that counts up while
the line is 1 and down while it is 0, stopping at 3 and 0. \cc{rx_bit} only
changes when the counter reaches 3 or 0, that is, when the line had the new
value for three samples in a row. Short glitches are ignored.

\svfile[style=numbered,firstline=92,firstnumber=92,lastline=125]{06_io_mirror/uart_rx.sv}

\textbf{Finding the middle of a bit.} \cc{os_cnt} has \cc{l2o-1} bits;
\cc{log2(8)} counts the binary digits of 8, which gives 4, so \cc{os_cnt} has
3 bits and counts the oversampling ticks modulo 8. In \cc{IDLE} it is held at
0: \cc{c ? a : b} is the \term{conditional operator}, a 2:1 multiplexer that
selects \cc{a} if \cc{c} is 1 and \cc{b} otherwise. When the filtered line falls to 0, the start bit has begun, and the state
machine goes to \cc{BIT_START}. Now \cc{os_cnt} counts, and \cc{sample_now}
is 1 at the tick on which \cc{os_cnt} is 3: about half a bit after the falling
edge, and from then on every 8 ticks, one bit later. Each \cc{sample_now}
advances the state, and in the eight data states (bit 3 of the state set) the
filtered bit is shifted into \cc{rx_data} from the left. After eight shifts,
the first bit received, $d_0$, has arrived at position 0. The delay of the
synchroniser and the filter does not matter: it delays the falling edge and
all data bits alike.

\svpart{style=numbered,firstline=128,firstnumber=128,lastline=132}{06_io_mirror/uart_rx.sv}

At the sample point of the stop bit, \cc{rx_ready} becomes 1 for one cycle if
the stop bit is 1; a frame without a valid stop bit is dropped. Since
\cc{sample_now} is 1 for one cycle, so is \cc{rx_ready}, as the comment says.
The rest of \cc{uart_rx.sv} measures pauses between bytes for
\cc{rx_idle} and \cc{rx_eop}, which we do not use.

\section{A Latch from an Incomplete case}
\label{sec:step06-latch}

An earlier version of \cc{test.sv} had no \cc{default} branch in the
\cc{always_comb} block of the state machine. With this version, yosys stops:

\outputfile{step06_latch.txt}

What is a latch doing here? An \cc{always_comb} block must say, for every
possible value of its inputs, what its outputs are. If there is an input value
for which \cc{next_state} is not assigned, then \cc{next_state} must keep its
previous value for that input. A circuit without clock that keeps a value is a
\term{latch}: level-sensitive storage, the D latch of the course. On an FPGA
it has to be built from LUTs with feedback, its timing is hard to analyse, and
it is almost never what we wanted. \cc{synth_ice40} therefore aborts by
default (its option \cc{-latches} has the default \cc{error}); with
\cc{-latches warn} it would only warn and build the latch from LUTs.

But the \cc{case} lists all four states. Why is it incomplete? Because
\cc{state} has 32 bits, as we saw: the \cc{case} covers four of its $2^{32}$
values. That only these four occur is a property of our state machine, not of
the circuit that computes \cc{next_state}, and yosys builds that circuit for
all inputs. (Verilator, which knows that an \cc{enum} variable holds one of its
values, does not complain. Simulator and synthesis tool look at the same text
differently.) The \cc{default} branch makes the assignment complete.

\begin{pitfall}
In \cc{always_comb}, assign every output in every case, for example with a
\cc{default} branch, or with a value assigned before the \cc{case}, as
\cc{dev_hex} does with \cc{digit = 0;}. In \cc{always_ff}, a missing
assignment is no problem: a flip-flop that keeps its value is a flip-flop with
enable. That is why the \cc{case} blocks in \cc{always_ff} above need no
\cc{default}.
\end{pitfall}

\begin{exercise}
\label{ex:step06-enum}
Remove the \cc{default} branch again, and give the \cc{enum} a base
type of two bits: write \cc{typedef enum logic [1:0]} instead of
\cc{typedef enum} in line 26. Does yosys still find a latch? Why?
\end{exercise}

\section{What yosys and nextpnr Make of It}

\outputfile{step06_yosys_stat.txt}

Most of the 40 \cc{SB_CARRY} belong to the two accumulators of the tick
generators: 18 each, for adders of 19 bits. The flip-flops come in more kinds
than before; the letters S and R stand for set and reset inputs, which yosys
uses for assignments of constants, such as \cc{acc <= inc} when
\cc{enable} is 0. nextpnr places it all on 133 logic cells:

\outputfile{step06_pnr.txt}

\section{Talking to the Board}
\label{sec:step06-kermit}

After \cc{make prog}, the board waits for characters. On the computer, the
USB chip appears as a serial port, and any terminal program can use it. The
repository contains a script for C-Kermit:

\par\noindent\begin{minipage}{\linewidth}
\lstinputlisting[style=term,title={\texttt{06\_io\_mirror/kermit-serial}}]{../06_io_mirror/kermit-serial}
\end{minipage}\par

\cc{set line} names the serial port. The USB chip has two channels, one for
programming the flash and one for the serial line. On a Mac the serial line
appears as \cc{/dev/tty.usbserial-} followed by an identifier of the chip; the
name in the script is that of the author's board. Look with
\cc{ls /dev/tty.usbserial-*} which name yours has, and change the line. (On Linux the port is usually
\cc{/dev/ttyUSB1}.) \cc{set speed 9600} is our baud rate, \cc{set stop-bits
2} matches the two stop bits of our transmitter. \cc{carrier-watch off} tells
Kermit not to wait for modem signals, which our line does not have, and the
two \cc{set terminal} lines make the Return key, which sends a carriage return
(code \cc{0D}), also start a new line on the screen. \cc{connect} starts the
terminal; \texttt{Ctrl-\textbackslash} followed by \cc{c} returns to the Kermit
prompt, where \cc{exit} ends it.

Copy the script to a directory in your \cc{PATH}, make it executable and run
it. Each key you press is sent to the FPGA, comes back and appears on the
screen, and the display shows its code: \cc{41} for \cc{A}, \cc{61} for
\cc{a}. Kermit itself does not show what you type, so if nothing appears,
the FPGA did not answer.

The design has a buffer of only one byte. While it waits in
\cc{IO_BUFFERING_DATA} for the transmitter, it ignores \cc{rx_ready}, and a
byte that arrives in this time is lost. When you type, this does not happen;
when a program sends many bytes in a row, it can. The next step puts FIFOs
between the UART and the rest of the design.

\section{Simulation}
\label{sec:step06-sim}

To see the UART at work, the testbench has to play the computer: it sends a
frame on \cc{RX} and watches what comes back on \cc{TX}. The level of
\cc{RX} in each clock cycle is a function of the cycle number:

\cpppart{firstline=9,lastline=29}{06_io_mirror/tb_test.cpp}

The main loop drives the clock itself, one cycle per iteration, prints every
change of \cc{TX} with its cycle, and samples the echo in the middle of each
data bit:

\cpppart{firstline=46,lastline=74}{06_io_mirror/tb_test.cpp}

The output of \cc{make sim}:

\outputfile{step06_echo.txt}

The middle of the stop bit of the frame on \cc{RX} is at cycle
$1000 + 9.5 \cdot 1250 = 12\,875$. The receiver reports the byte somewhat
later, since the synchroniser and the filter delay the line by a few
oversampling ticks, and the echo starts at cycle 13\,757 with the start bit.
Then, with $\mathtt{41}_{16}$ = $\mathtt{01000001}_2$ sent LSB first: $d_0 =
1$ from 15\,009, $d_1$ to $d_5 = 0$ from 16\,260, $d_6 = 1$ from 22\,516, $d_7
= 0$ from 23\,768, and from 25\,019 on the line stays 1 for the stop bits and
idle. The bits last 1251 or 1252 cycles, as computed in
Section~\ref{sec:step06-baud}; five bits from 16\,260 to 22\,516 are 6256
cycles. The testbench decodes the echo and compares it with the byte it sent;
if they differ, it prints \cc{FAILED} and \cc{make sim} ends with an error.

\begin{exercise}
\label{ex:step06-tb}
Change \cc{rx_level} so that the testbench sends the 26 letters \cc{A} to
\cc{Z} back to back, each frame directly after the stop bit of the one
before, and print the bytes that come back. Which letters are missing, and
why? (Hint: compare the length of a frame on \cc{RX} and on \cc{TX}.)
\end{exercise}

\section{What Is Next}

We can now exchange bytes with the computer, one at a time. Step 07 turns this
into an I/O interface for the rest of the design: FIFOs that buffer bytes in
both directions, built from block RAM, and the first SystemVerilog
\cc{interface}s to connect the parts.

\solutions

\solution
In \cc{IO_INIT_TRANSMISSION}, send the next code instead of the received one:
\begin{lstlisting}
tx_data_in <= data_buf + 1;
\end{lstlisting}
The display still shows the received byte, since it shows \cc{data_buf}. (For
\cc{z}, the answer is \texttt{\{}, code \cc{7B}.)

\solution
With $f = 12\,000\,000$, $b = 115\,200$, $o = 1$: $f / b = 104$ (integer
division), which has 7 binary digits, so \cc{acc_width} $= 15$. Then
\cc{shiftlimiter} $=$ \cc{log2_(115200 >> 16)} $=$ \cc{log2_(1)} $= 1$, and
\[
  \mathit{inc} = \left\lfloor \frac{115\,200 \cdot 2^{14} + 3\,000\,000}{6\,000\,000}
  \right\rfloor = \lfloor 315.07 \rfloor = 315 .
\]
The simulation with \cc{$display} agrees (\cc{acc_width 15, shiftlimit 1, inc
315}). The tick rate is $12\,000\,000 \cdot 315 / 2^{15} \approx 115\,356$ per
second, $0.14\,\%$ too fast. Without rounding, $115\,200 \cdot 2^{15} / 12\,000\,000
= 314.57$; the error comes from rounding to an integer, and a larger $w$ would
make it smaller.

\solution
With \cc{logic [1:0]} as base type, \cc{state} has 2 bits, and the four
branches cover all four values. The \cc{case} is complete, \cc{next_state} is
assigned for every input, and there is no latch. yosys synthesises the design
without error, and the statistics are the same as with the \cc{default}
branch.

\solution
In \cc{rx_level}, the frame number is $f = \lfloor (c - \mathtt{T0}) /
(10 \cdot \mathtt{BIT}) \rfloor$ and the bit within the frame is $k =
\lfloor (c - \mathtt{T0}) / \mathtt{BIT} \rfloor \bmod 10$; for $f < 26$
the function returns the start bit, a data bit of letter number $f$, or the
stop bit, as before, and 1 after the last frame. The
echo then misses \cc{L} and \cc{W}: 24 letters come back. A frame on \cc{RX}
takes 10 bit times, 12\,500 cycles, but our transmitter sends two stop bits,
11 bit times or about 13\,768 cycles. So every character, the echo falls
behind by about 1268 cycles. The state machine accepts a byte only in
\cc{IO_READY_TO_RECEIVE}; with a byte waiting in \cc{IO_BUFFERING_DATA} for
the transmitter, it ignores \cc{rx_ready}. After ten characters the backlog
is a whole frame: the eleventh letter, \cc{L}, arrives while \cc{K} is still
waiting, and is lost. Then the backlog starts again from a smaller value, and
eleven letters later \cc{W} is lost. The FIFOs of step 07 make the buffer
much longer, but the rate problem remains: in the long run a sender that is
faster than the transmitter fills any buffer.
